\chapter*{Введение}
\addcontentsline{toc}{chapter}{Введение}
\label{ch:intro}

Современная веб-страница --- это не статический документ, а полноценное
клиентское приложение. Пользователь ожидает, что интерфейс реагирует на
действия мгновенно, введённые данные не теряются при перезагрузке, а
результат работы можно сохранить в привычном формате. Все эти возможности
реализуются тремя базовыми технологиями браузера: \textbf{HTML} описывает
структуру, \textbf{CSS} --- оформление и анимации, \textbf{JavaScript} ---
поведение и работу с данными.

Предыдущая лабораторная работа (\textnumero\,13) была посвящена основам
JavaScript и построению сайта на данных внешнего API. Текущая работа
развивает практические навыки клиентской разработки на примере
\textit{интерактивного резюме}: страница не просто отображает информацию,
а позволяет редактировать её прямо в браузере, сохраняет правки между
сессиями и экспортирует готовый документ в формат PDF.

В качестве темы выбрана веб-страница с резюме, свёрстанная по типовому
шаблону. Каждый текстовый элемент резюме редактируется <<на месте>>
(атрибут \texttt{contenteditable}), все изменения сохраняются в хранилище
браузера \texttt{localStorage} и восстанавливаются при следующем открытии.
Кнопка <<Скачать PDF>> формирует документ средствами самого браузера. Любое
изменение элемента сопровождается CSS-анимацией, а клики по элементам ---
эффектом \textbf{Material Wave} (<<волна>>, ripple) из дизайн-системы
Google Material Design. Готовое приложение публикуется на собственном
домашнем сервере под обратным прокси \textbf{Caddy} и доступно по адресу
\url{https://resume.server34.netcraze.club}.

Согласно требованиям методички, разработка ведётся \textbf{только базовыми
технологиями} --- HTML, CSS и JavaScript --- без сторонних библиотек и
фреймворков и без этапа сборки.

Цель работы --- получить практические навыки разработки интерактивных
веб-приложений: реализовать редактирование контента, сохранение данных на
стороне клиента, экспорт страницы в PDF, CSS-анимации и адаптивную вёрстку.

Задачи:
\begin{enumerate}
  \item Сверстать веб-страницу с резюме по типовому шаблону, используя
        семантическую разметку HTML5.
  \item Реализовать редактирование текстовых элементов резюме с сохранением
        изменений внутри самих элементов.
  \item Добавить кнопку <<Скачать>>, преобразующую текущую версию резюме в
        формат PDF и начинающую загрузку файла.
  \item Анимировать изменения элементов средствами CSS-анимаций.
  \item Обеспечить адаптивность страницы и эффект Material Wave при кликах.
  \item Реализовать сохранение введённых данных при перезагрузке страницы и
        опубликовать приложение на публичном хостинге.
\end{enumerate}

\textbf{Стенд:}
\begin{itemize}
  \item современный браузер (Chrome, Firefox, Safari) --- среда исполнения
        приложения; страница открывается как напрямую через \texttt{file://},
        так и с веб-сервера;
  \item редактор кода с линтерами \texttt{JSHint}\,/\,\texttt{ESLint} ---
        проверка качества JavaScript-кода;
  \item валидатор разметки \texttt{validator.w3.org} --- проверка семантики
        вёрстки;
  \item собственный домашний сервер (Debian~12, Docker Compose) с обратным
        прокси \texttt{Caddy} и автоматическим выпуском TLS-сертификатов
        Let's~Encrypt --- площадка для публикации готового сайта.
\end{itemize}

\endinput
